% !TeX root = ../papers/nc-p3.tex
\section{種数三の切断グラフと二重被覆・正規四次の排除}
\label{sec:sectional-eighteen-kernel-graph}
\begin{theorem}\label{thm:sectional-eighteen-kernel-graph}
定理\ref{thm:genus-three-kernel-excess}の十八次の場合、滑らかな二次核曲面への
次数二の核写像は存在しない。また四次核曲面で入力逆写像が余次元一において正則となる
場合は生じず、特に正規四次核曲面はない。
\end{theorem}
\begin{proof}[平面上の有限グラフ]
$G$ は一般入力平面八次 $E$ と核空間の積への正規化曲線 $B$ の像スキームである。
実際定理\ref{thm:sectional-cofactor-length}の等号により $M_E$ は $G$ 上の可逆層 $P$ の押出し。
入力の各点上の有限半局所環で、四つの核座標の一般結合は全ての剰余点で単元となって
$P$ を自明化する。これで割った座標は押出し加群の生成元だから
$\mathcal O_E$ 上の代数としても $\mathcal O_G$ を生成する。
従って $(h,j):G\to E\times\mathbb P^3$ は閉埋込みである。
核像 $J=j(G)$ への射は双有理である。核曲面の一般点での対の次元と、
異なる二入力点を含む平面の一次元族を数えれば、一般平面上で同じ核を持つ対は
その開部分では高々有限個となるためである。
$J$ は非平面である。二次曲面または三次・四次曲面の平面切断が整六次曲線を含むことはない。
よって $3\le p_a(J)\le4$。三なら $J=G$、四なら $J$ は二次と三次の完全交叉。
どちらもGorensteinである。

四次の場合、$8\deg(K)\deg(\text{核写像})\le36$ より核写像は双有理である。
$K$ の双対化層は自明なので、$J$ のイデアルの双対の列
$0\to\mathcal O_K\to\mathcal Hom_K(I_J,\mathcal O_K)\to\omega_J\to0$
を使える。逆写像が余次元一で正則なら、一般入力平面の切断で割った四座標の比は
この中項の切断となる。定理\ref{thm:cofactor-eighteen-normal-quartic}の随伴の議論の通り、
有限集合を越える延長は $S_2$ 性により、一般点で右項への像の比は入力平面座標を与える。
しかし $\deg\omega_J\le6$ で、入力平面での次数八には達しない。
\end{proof}

\begin{proof}[滑らかな二次曲面への二対一写像]
核二母線の引戻しの入力次数は $(2,4)$ または $(3,3)$ である。
定理\ref{thm:cofactor-eighteen-no-smooth-quadric}の被覆の数値計算は、
$H^2=8$, $L^2=4$, $HL=6$ と一般入力切断の正規化種数三以下を使う。
同定理のうち階数二曲線の不可約性を必要とした跡の零点での議論を、
上のグラフの閉埋込みで置き換える。

$(2,4)$ ではHodge等号と有限入力跡から正規化は
$S=\mathbb P^1_u\times\mathbb P^1_v$、$H=(2,2)$、核射は $(u,v)\mapsto(a(u),v)$、
$\deg a=2$。跡の零因子は類 $(2,0)$ で、有限個の $u$ 上にある。
一般入力平面の正規化 $B$ は種数一の $(2,2)$ 曲線である。
跡因子に属さない全ての枝では、跡包含の両端がその枝で一致するので $M_E$
は正規化上の値をそこで独立に与える。特に $G$ の正規化写像が同型でない可能性は
跡因子と一般平面との交点に限られる。
一般 $v$ の近くで $(f,v)$ は閉埋込みである。
実際 $v$ 固定の母線を入力次数二で双有理に写すので、その三つの二次切断が
完全系を張ることは開条件であり、相対二次埋込みが得られる。
一般平面で跡因子上の点はこの開部分にある（各該当母線の入力は非定数）。
従ってその点でも、より多くの座標を持つ $(f,a(u),v)$ が全ての枝と接方向を分離する。
ゆえにグラフは全点で $B$ に等しい。$p_a(G)=3$ と $g(B)=1$ に反する。

$(3,3)$ では同定理の随伴と線織面計算により、正規二重被覆 $T$ は $(2,2)$ 分岐で
有理二重点のみを持つ。最小解消 $S$ 上で $K_S=-L$、入力束 $H=L+P_1$ は自由で
$P_1$ は $HP_1=LP_1=2$ の自由鉛筆。一般入力切断 $B$ は種数二である。
入力に潰れる曲線は全て $L$ 次数零であり、有限な開部分での有効跡因子は
$(3L-2H)L=0$ のため $L$ 次数零の曲線にしか現れない。
核被覆の各特異点の像 $y$ を通る核平面の固定引戻し部分 $Z_y$ に対しては
$-2\le Z_y^2<0$、$H-3L/2$ の自乗が $-1$ なので $HZ_y\le\sqrt{2}$、
従って $HZ_y\le1$ である。ここまではいずれも被覆と交点数の計算である。
$B$ がこの部分に会うならただ一点で横断し、ある核平面の式はそこで単純零点を持つ。
また特異被覆点は分岐上にあり、その $y$ 上に別の被覆点はない。
従ってこの点で $B\to J$ は一つの逆像を持ち非分岐で、局所的に閉埋込みである
（有限射の繊維長一でNakayamaの補題を用いる）。他の全ての点では跡包含が
正規化の値を与えるので、ここでもグラフ全体は $B$ に一致する。
$g(B)=2$ と $p_a(G)=3$ が矛盾する。
\end{proof}
\path{computations/verify_sectional_eighteen_kernel_graph.py} は被覆の交点数、跡因子、
固定例外部分の次数と比較に使う三つの種数を検算する。
